\documentclass[11pt]{article}
\usepackage[margin=1in]{geometry}
\usepackage{amsmath,amssymb,amsthm}
\usepackage[hidelinks]{hyperref}
\newtheorem{proposition}{Proposition}
\newtheorem{lemma}{Lemma}
\title{Strong avoidance does not force local principal ideal domains}
\author{Principia Math in Progress\\MathDB 353853}
\date{12 September 2026}
\begin{document}
\maketitle
\begin{abstract}
Every commutative ring with linearly ordered ideals has strong avoidance in the
sense of Chen and Tarizadeh. This gives a four-element counterexample to their
Conjecture 4.10. An explicit valuation domain with a nonprincipal maximal ideal
also disproves the conjecture within integral domains.
\end{abstract}

\paragraph{Task 1. Resolve the strong-avoidance characterization.}
We disprove the characterization in Conjecture 4.10 of \cite{CT}, using its
actual Definition 4.2. All rings here are commutative with identity $1\ne0$,
and ring maps preserve the identity. For a ring map $\varphi:R\to S$, write
$IS$ for the ideal generated by $\varphi(I)$. The ring $S$ has
\emph{strong avoidance} when, for every such map and every finite nonempty
family of ideals $I,I_1,\ldots,I_n$ of $R$,
\[
 I\subseteq\bigcup_{j=1}^n I_j
 \quad\Longrightarrow\quad
 IS\subseteq I_jS\text{ for some }j.
\]
The conjecture asserts that this is equivalent to $S_{\mathfrak p}$ being
a principal ideal domain (PID) for every prime ideal $\mathfrak p$ of $S$.
No Noetherian or reduced hypothesis appears in the source.

The MathDB paraphrase instead describes a condition on two ideals whose
union is the whole ring. That condition is automatic: one of those ideals
contains $1$. Our counterexamples satisfy the stronger, actual source definition.

\begin{lemma}\label{chain}
If the ideals of $S$ are linearly ordered by inclusion, then $S$ has strong
avoidance.
\end{lemma}
\begin{proof}
Fix an arbitrary map $\varphi:R\to S$ and a cover as above. Among the finite
family of ideals $I_1S,\ldots,I_nS$, choose a largest one, say $I_jS$. Then
\[
 \varphi(I)\subseteq\bigcup_{k=1}^n\varphi(I_k)
 \subseteq\bigcup_{k=1}^n I_kS=I_jS.
\]
Since $I_jS$ is an ideal, it contains the ideal generated by $\varphi(I)$,
namely $IS$. The map need not be injective or surjective, and $I$ need not
be finitely generated.
\end{proof}

\begin{proposition}\label{finite}
The ring $S=\mathbb Z/4\mathbb Z$ has strong avoidance, but its localization
at its unique prime ideal is not a PID.
\end{proposition}
\begin{proof}
Its ideals are $(0)$, $(2)$ and $S$, so Lemma~\ref{chain} applies.
Any prime ideal contains the nilpotent element $2$; the quotient by $(2)$ is
$\mathbb F_2$, so $(2)$ is the unique prime ideal. Every element outside
$(2)$ is a unit: $1$ and $3$ are their own inverses. Thus $S_{(2)}=S$.
The element $2$ is nonzero and has square zero, so $S$ is not an integral
domain and therefore is not a PID.
\end{proof}

The same argument applies to $\mathbb F_2[\epsilon]/(\epsilon^2)$, whose
ideals are $(0)$, $(\epsilon)$ and the whole ring. The next example shows
that the obstruction persists even for domains, and even if ``PID'' were
replaced by ``principal ideal ring.''

\begin{proposition}\label{domain}
There is a valuation domain $V$ with strong avoidance whose maximal ideal
$\mathfrak m$ is not principal. In particular, $V_{\mathfrak m}=V$ is not a PID.
\end{proposition}
\begin{proof}
Let $t_0$ be transcendental over $\mathbb F_2$, and in an algebraic closure
of $\mathbb F_2(t_0)$ choose elements $t_n$ with $t_{n+1}^2=t_n$ for $n\ge0$.
Set
\[
 K=\bigcup_{n\ge0}\mathbb F_2(t_n),\qquad
 V=\bigcup_{n\ge0}\mathbb F_2[t_n]_{(t_n)}.
\]
Both unions are nested. For the second, a polynomial denominator with
nonzero constant term retains that property after substituting
$t_n=t_{n+1}^2$. In particular, $K$ is a field and $V$ is a subring of $K$.

On $\mathbb F_2(t_n)^\times$ define
\[
 v(f)=2^{-n}\operatorname{ord}_{t_n}(f).
\]
Substitution doubles the order of every nonzero rational function, so these
definitions agree on overlaps and define a valuation
$v:K^\times\to\mathbb Z[1/2]$. Every dyadic rational occurs, since
$v(t_n^a)=a/2^n$ for $a\in\mathbb Z$. At each finite stage, nonnegative
order is exactly the condition for membership in its local polynomial ring.
Consequently
\[
 V=\{0\}\cup\{f\in K^\times:v(f)\ge0\}.
\]
For every nonzero $f\in K$, either $f$ or $f^{-1}$ belongs to $V$.

To check that all ideals of $V$ are comparable, suppose $A\nsubseteq B$
and choose $a\in A\setminus B$. For nonzero $b\in B$, the element $a/b$
cannot belong to $V$, since that would imply $a\in B$. Hence $b/a\in V$
and $b\in A$. Thus $B\subseteq A$. Lemma~\ref{chain} now proves strong
avoidance for $V$.

The nonunits of $V$ form the ideal
\[
 \mathfrak m=\{0\}\cup\{f\in K^\times:v(f)>0\}:
\]
the valuation inequality shows closure under addition, and multiplication
by elements of $V$ preserves positivity. Elements of value zero have
their inverses in $V$. Thus $\mathfrak m$ is the unique maximal ideal,
and localizing there changes nothing.

If $\mathfrak m=(a)$, then $a\ne0$ and $q=v(a)>0$. Write $q=c/2^n$ with
integers $c>0$ and $n\ge0$. The element $b=t_{n+1}^c$ has value $q/2>0$,
so $b\in\mathfrak m$. But $v(b/a)=-q/2<0$, whence $b/a\notin V$ and
$b\notin(a)$, a contradiction. Therefore $\mathfrak m$ is not principal.
\end{proof}

Either proposition disproves the necessary direction of the stated
equivalence. Proposition~\ref{domain} does so within integral domains.
These counterexamples do not assert a complete replacement classification.

\paragraph{Task 2. Classify finite local examples.}
Determine which finite commutative local rings have strong avoidance.
In particular, decide whether every such ring must have its ideals linearly
ordered, and identify the role of the residue field and nilpotent ideals.

\paragraph{Task 3. Determine the effect of finite products.}
Find necessary and sufficient conditions for a finite product of rings
with strong avoidance to retain that property. Track when the containing
ideal can be chosen with the same index in every factor.

\paragraph{Task 4. Investigate the domain classification.}
Determine which Pr\"ufer domains have strong avoidance, and whether every
domain with strong avoidance must be Pr\"ufer. The valuation-domain lemma
settles the local sufficiency direction but leaves the global choice of
an index unresolved.

\paragraph{Task 5. Identify valid preservation statements.}
Determine when strong avoidance passes to localizations, quotients and
directed unions. Use those results to formulate and test a corrected
characterization under explicit finiteness or local hypotheses.

\begin{thebibliography}{9}
\bibitem{CT}
J. Chen and A. Tarizadeh,
\emph{On the ideal avoidance property},
\href{https://arxiv.org/abs/2203.04256v4}{arXiv:2203.04256v4}
(6 September 2023), Definition 4.2, p.~12, and Conjecture 4.10, p.~14.
Journal reference: J. Pure Appl. Algebra \textbf{228}(3) (2024), 107500,
\href{https://doi.org/10.1016/j.jpaa.2023.107500}{doi:10.1016/j.jpaa.2023.107500}.
The proof uses the current author-linked arXiv version; the final typeset
journal text was unavailable during the audit.
\end{thebibliography}
\end{document}
